\begin{python0}
from solutions import *; clear()
\end{python0}
\sectionthree{When do I use the \texttt{for}-Loop?}

I've already shown you that you can always rewrite a \texttt{for}-loop into a \texttt{while}-loop and vice versa.

So ... how do you choose which type of loop to use?

The general advice is simple: If your loop has a control (index or
counter) variable that runs across a sequence of values without user
intervention and the control variable changes in a predictable way, you
should use the for-loop. The user might need to ``configure'' this
sequence: for instance, by specifying a starting value and an ending
value and how to step through the intermediate values. But once
that's done, the control variable in the for-loop should
run without user intervention.

Otherwise you should use the \texttt{while}-loop.

Look at this familiar example:
\begin{console}
int start = 0, end = 0, step = 0;
std::cin >> start >> end >> step;

int sum = 0;
for (int i = start; i <= end; i += step)
{   
    sum += i;
}
std::cout << sum << std::endl; 
\end{console}

Note that the loop is controlled by \texttt{i} and \texttt{i} steps through
a sequence with the starting and ending values known before the loop
starts.

Note that the control variable is only used in loop through some
statements to achieve a goal. Look at the above example: The control
variable is defined within the for-loop. After the for-loop, the control
variable is destroyed.

The next section shows you a situation where the data to process is
entered by the user and hence cannot be predicted ahead of time.

